\documentclass[a4paper]{article}
\usepackage[T1]{fontenc}
\usepackage{CJKutf8}
\usepackage[margin=1.2cm]{geometry}
\usepackage{tikz}
\usetikzlibrary{shapes.geometric, arrows.meta}

\begin{document}
\begin{CJK}{UTF8}{gbsn}
\thispagestyle{empty}
\centering

\begin{tikzpicture}[
  font=\small,
  % --- node style definitions ---
  process/.style={rectangle, draw=black!70, thick, rounded corners=1mm, fill=blue!6,
                  text width=47mm, align=center, minimum height=10mm, inner sep=2.2mm},
  decision/.style={diamond, draw=black!70, thick, fill=yellow!15, aspect=1.8,
                   text width=20mm, align=center, inner sep=0.5mm},
  terminator/.style={rectangle, rounded corners=3mm, draw=black!70, thick, fill=green!12,
                      minimum width=20mm, minimum height=8mm, align=center},
  lbl/.style={font=\scriptsize, fill=white, inner sep=1pt},
  arr/.style={-{Stealth[length=2.2mm]}, thick},
]

% ================= main spine (top to bottom) =================
\node[terminator]                 (start) at (0, 0)     {开始};
\node[process]                    (s1)    at (0, -1.7)  {步骤1：将修改6个宏的固件\\发给施耐德Pentest团队测试};
\node[decision]                   (d1)    at (0, -3.9)  {测试\\通过？};
\node[process]                    (s2)    at (0, -6.3)  {步骤2：将瑞萨标准硬件平台+\\标准固件发给施耐德\\Pentest团队测试};
\node[decision]                   (d2)    at (0, -9.0)  {测试\\通过？};

% ================= branch after d2 (side by side) =================
\node[process]                    (s4)    at (-4.6, -12.3) {步骤4：固高将固件代码发给瑞萨，\\颜工与FAE杨工对比标准固件，\\共同排查导致Pentest测试\\不通过的改动点};
\node[process]                    (s5)    at (4.6, -11.7)  {步骤5：FAE杨工负责解决标准\\固件中导致Pentest测试\\不通过的问题};
\node[process]                    (s6)    at (4.6, -14.4)  {步骤6：颜工负责将FAE杨工的\\改动点合并到固高固件中};

% ================= merge point =================
\node[process]                    (s7)    at (0, -17.0) {步骤7：固高固件再次进行\\Pentest测试，验证问题已解决};
\node[terminator]                 (end)   at (0, -19.0) {结束};

% ================= side terminator for the 1a "pass" exit =================
\node[terminator]                 (done)  at (5.0, -3.9) {完成};

% ================= arrows =================
\draw[arr] (start) -- (s1);
\draw[arr] (s1) -- (d1);
\draw[arr] (d1) -- node[lbl, above] {是} (done);
\draw[arr] (d1) -- node[lbl, left]  {否} (s2);
\draw[arr] (s2) -- (d2);

% d2 splits left (pass -> s4) and right (fail -> s5)
\draw[arr] (d2.west) -| node[lbl, pos=0.22, above] {是} (s4.north);
\draw[arr] (d2.east) -| node[lbl, pos=0.22, above] {否} (s5.north);

\draw[arr] (s5) -- (s6);

% both branches converge into s7
\draw[arr] (s4.south) |- (s7.west);
\draw[arr] (s6.south) |- (s7.east);

\draw[arr] (s7) -- (end);

\end{tikzpicture}

\end{CJK}
\end{document}
